\section{Fiber ellipsoids and joint positivity}\label{sec:ellipsoids}

We turn the disc inequality into a scalar comparison by optimizing a
Hermitian form in each cotangent fiber.  Throughout this Section, $t$ is
physical time, and $0<t<T$ for a fixed finite $T$.  Thus $t$ plays the
role of the center-time parameter in Proposition~\ref{prop:disc-inequality}.
Let $v$ be its lower semicontinuous disc envelope.  We use
\begin{equation}\label{ell:envelope-bounds}
 N(x,\xi,0)\le v(x,\xi,t)\le N(x,\xi,t),
 \qquad v(x,c\xi,t)=|c|^2v(x,\xi,t).
\end{equation}
The matrices below represent Hermitian forms on cotangent columns by
$\xi^*B\xi$.  Accordingly, contraction against a horizontal Levi form
$H$ is $\tr(B^{\mathsf T}H)$; in particular,
$\Delta_{h_t}f=\tr(b_t^{\mathsf T}\dd_xf)$.

\subsection{The optimizer and its contacts}

The optimizer/contact-resolution argument is a Hermitian adaptation of
John's optimal-ellipsoid contact method \cite{John1948}. The
lower-semicontinuity details are proved here; the common-base product
support and the ensuing PDE comparison are further arguments below,
not consequences imported from the classical real convex-body theorem.

For each $(x,t)$, let $B(x,t)$ be the feasible matrix maximizing the
determinant:
\begin{equation}\label{ell:optimizer}
 \det B(x,t)=\max\bigl\{\det D:
 D>0,\ \xi^*D\xi\le v(x,\xi,t)\text{ for every }\xi\bigr\},
 \qquad \mathcal L=\log\det B.
\end{equation}
The ellipsoid $\{\xi:\xi^*D\xi<1\}$ of a feasible matrix contains
$\{\xi:v(x,\xi,t)<1\}$.  Maximizing $\det D$ minimizes its Euclidean
volume in these fiber coordinates.

\begin{lemma}\label{ell:optimizer-contacts}
The maximum in \eqref{ell:optimizer} exists and is unique.
The function $\mathcal L$ is lower semicontinuous in each holomorphic
coordinate chart, and
\begin{equation}\label{ell:volume-gap}
 F_0:=\log\det b_t-\mathcal L
 \quad\text{satisfies}\quad 0\le F_0\le\rho.
\end{equation}
At any fixed point, choose linear cotangent coordinates in which $B=I$.
There are finitely many unit columns $\xi_i$ and positive numbers
$\alpha_i$ such that
\begin{equation}\label{ell:contact-resolution}
 v(x,\xi_i,t)=1,\qquad
 \sum_{i=1}^m\alpha_i\xi_i\xi_i^*=I,
 \qquad \sum_{i=1}^m\alpha_i=n.
\end{equation}
The largest weight can be made arbitrarily small by repeating columns.
\end{lemma}

\begin{proof}
The initial cometric $b_0$ is feasible by
\eqref{ell:envelope-bounds}.  Every feasible matrix satisfies $D\le b_t$.
The constraints therefore define a compact set if positive semidefinite
matrices are temporarily admitted.  A determinant maximizer in that set
has determinant at least $\det b_0>0$, so is positive definite.
Strict concavity of $\log\det$ on the positive cone gives uniqueness.
It also follows that $\det b_0\le\det B\le\det b_t$, which proves
\eqref{ell:volume-gap}, since
$\rho=\log\det b_t-\log\det b_0$.

Fix a feasible positive matrix $D$ at $(x_0,t_0)$.  For $0<\varepsilon<1$,
the candidate $(1-\varepsilon)D$ has a strictly positive margin on the
unit fiber sphere.  Lower semicontinuity of $v$ and compactness of that
sphere preserve this margin for all nearby $(x,t)$.  Consequently
\[
 \liminf_{(x,t)\to(x_0,t_0)}\mathcal L(x,t)
 \ge\log\det D+n\log(1-\varepsilon).
\]
Take $D=B(x_0,t_0)$ and then $\varepsilon\downarrow0$.
This proves the asserted semicontinuity without asserting continuity of
the optimizing matrix.  Under a change of holomorphic base coordinates,
both $\mathcal L$ and $\log\det b_t$ acquire the same time-independent
pluriharmonic summand.  Thus $F_0$ is a globally defined upper
semicontinuous function.

Normalize $B(x_0,t_0)=I$, using a constant linear base-coordinate change
and its induced cotangent change.  On the unit sphere put
\[
 E=\{\xi:v(x_0,\xi,t_0)=1\}.
\]
This is compact, because $v\ge|\xi|^2$ at the chosen base point.
Suppose that a Hermitian matrix $H$ is strictly negative on $E$.
It is uniformly negative in a neighborhood of $E$ on the sphere;
on the complementary compact set, $v-1$ has a positive minimum.
It follows that $I+\varepsilon H$ is feasible for all sufficiently
small $\varepsilon>0$.  Optimality then implies $\tr H\le0$.

The cone generated by $\xi\xi^*$, $\xi\in E$, is closed: its generators
have trace one, so bounded trace bounds the total coefficient in any
conic combination.  If $I$ were outside this cone, finite-dimensional
separation would give $H$ with $\tr H>0$ and $\xi^*H\xi\le0$ on $E$.
Subtracting a sufficiently small positive multiple of $I$ contradicts
the preceding paragraph.  Hence $I$ belongs to the cone.  A finite
conic representation gives \eqref{ell:contact-resolution}, after
discarding zero coefficients.  Taking its trace gives $\sum_i\alpha_i=n$.
Replacing each pair $(\xi_i,\alpha_i)$ by $r$ copies of
$(\xi_i,\alpha_i/r)$ preserves every identity and divides the largest
weight by $r$.
\end{proof}

\begin{proposition}\label{prop:ellipsoid-comparison}
For the optimizer in \eqref{ell:optimizer}, every smooth lower test
$\lambda$ of $\mathcal L$ at an interior point $(x_0,t_0)$ satisfies
\begin{equation}\label{ell:volume-viscosity}
 \partial_t\lambda(x_0,t_0)
 \ge e^{-F_0(x_0,t_0)}\Delta_{h_{t_0}}\lambda(x_0,t_0).
\end{equation}
Thus $\mathcal L$ satisfies this inequality in the lower viscosity
sense on $M\times(0,T)$.
\end{proposition}

To prove this scalar inequality, we construct simultaneous contact
jets at a common base point and time and transfer them to lower tests
of $v$.  The transferred tests must have positive vertical Levi blocks,
as required by
Proposition~\ref{prop:disc-inequality}.  We first construct a support
for the weighted sum of the contact values, then obtain simultaneous
contacts and add their Schur-complement inequalities.

\subsection{A product support with positive vertical Levi form}

Fix a smooth lower test $\lambda$ of $\mathcal L$ at $(x_0,t_0)$,
and choose the normalized contact list in that fiber.  Allow the
columns to vary independently, and set
\begin{equation}\label{ell:rectangular-matrix}
 M_c=(\sqrt{\alpha_1}\xi_1,\ldots,\sqrt{\alpha_m}\xi_m),
 \qquad A=M_cM_c^*.
\end{equation}
Near the original tuple $M_0$, the matrix $M_c$ has rank $n$.
Let $\Pi(M_c)$ be its row plane in $\Gr(n,m)$.
For every nearby $(x,t)$ and every full-rank tuple,
\begin{equation}\label{ell:product-support}
 \sum_i\alpha_i v(x,\xi_i,t)
 \ge\tr(B(x,t)A)
 \ge n+\mathcal L(x,t)+\log\det A.
\end{equation}
Indeed apply $a\ge1+\log a$ to the positive eigenvalues of
$B^{1/2}AB^{1/2}$.  Equality in the second inequality holds exactly
when $B^{1/2}AB^{1/2}=I$.

Here is the Levi calculation that will control each separate fiber.
For a full-rank complex $n\times m$ matrix $M$, put $A=MM^*$.
For a matrix variation $Z$, differentiating along $M+zZ$ gives
\begin{align}
 \dd\log\det(MM^*)[Z,\bar Z]
 &=\tr(A^{-1}ZZ^*)
   -\tr(A^{-1}ZM^*A^{-1}MZ^*) \notag\\
 &=\tr\bigl(A^{-1}Z(I-M^*A^{-1}M)Z^*\bigr)\ge0.
 \label{ell:rectangular-levi}
\end{align}
The middle factor in the final expression is an orthogonal projection.
This form is the pullback by $\Pi$ of the Grassmannian form induced
by the Pl\"ucker embedding.  In particular, for a variation $V_i$ in
the $i$th unweighted column alone,
\begin{equation}\label{ell:single-column}
 \dd_{\xi_i}\log\det A[V_i,\bar V_i]
 =\alpha_i\bigl(1-\alpha_i\xi_i^*A^{-1}\xi_i\bigr)
      V_i^*A^{-1}V_i.
\end{equation}
At any tuple with $A=I$ and $|\xi_i|=1$, the form after division by
$\alpha_i$ is therefore $(1-\alpha_i)I$.  Repeating the contact
columns leaves their matrix resolution unchanged and makes these
normalized vertical forms as close to $I$ as needed.

Subtracting a small multiple of
$(|x-x_0|^2+|t-t_0|^2)^2$ from $\lambda$ makes its contact strict on a compact
coordinate cylinder without changing its time derivative or spatial
Levi form there.  Choose a smooth nonnegative function $\eta$ on
$\Gr(n,m)$, vanishing only at $\Pi(M_0)$, and scale it so that
\begin{equation}\label{ell:grassmann-penalty}
 \dd(\eta\circ\Pi)\le\dd\log\det A.
\end{equation}
Such a function exists: the squared Euclidean distance from the
orthogonal projector of a plane to the projector of $\Pi(M_0)$ is
smooth and vanishes at that plane alone, and its Levi form is bounded
above by a constant times the positive Grassmannian metric.
Fix $0<\theta<1$, and use the smooth support
\begin{equation}\label{ell:localized-support}
 g_\theta(x,t,\xi_1,\ldots,\xi_m)
 =n+\lambda(x,t)+\log\det A-\theta\eta(\Pi).
\end{equation}
The difference between the left side of \eqref{ell:product-support}
and $g_\theta$ is nonnegative.  At equality, strictness forces
$(x,t)=(x_0,t_0)$, the trace inequality forces $A=I$, and the penalty
forces $\Pi=\Pi(M_0)$.  Thus $M_c=UM_0$ for a unitary $U\in U(n)$.
In particular each unweighted column still has norm one.  Equality
also requires each column to remain a contact of $v$.
The equality set is therefore a compact \emph{subset} of this unitary
orbit.  It is nonempty, since it contains the original contact tuple.

Choose a bounded open neighborhood $O$ containing this equality set,
with compact closure in the base-time cylinder and in the full-rank
matrix locus.  Its boundary avoids the equality set.  Lower
semicontinuity gives a positive gap on $\partial O$ between the sum
and its support.  All columns on $\overline O$, and their small
neighborhoods, lie in one bounded cotangent-coordinate region.

\subsection{Simultaneous contacts with a common base}

We use Alexandrov's almost-everywhere second-order differentiability
theorem \cite{Aleksandrov1939}, in the form of
\cite[Appendix A, Theorem~A.2]{CIL}.  The contact-set and translated-test
arguments below follow the Jensen and inf-convolution framework
of \cite[Appendix A, Lemmas~A.3--A.5]{CIL}.  We give the local details
that retain one shared base and time and introduce no negative
error in the vertical Levi form.

Use the real coordinate vector $q=(\Re x,\Im x,t,\Re\xi,\Im\xi)$ and
define, with sources restricted to a fixed slightly larger compact
coordinate region,
\begin{equation}\label{ell:inf-convolution}
 v^\kappa(q)=\inf_z\left\{v(z)+\frac{|q-z|^2}{2\kappa}\right\}.
\end{equation}
The infimum is attained.  On the region used here,
$v^\kappa\le v$ and any minimizing source satisfies
$|z-q|\le C\sqrt\kappa$, by the local upper bound and nonnegativity
in \eqref{ell:envelope-bounds}.  Sources are consequently interior
for small $\kappa$.  Moreover,
\begin{equation}\label{ell:epigraph-limit}
 q_j\to q,\quad\kappa_j\downarrow0
 \quad\Longrightarrow\quad
 \liminf_j v^{\kappa_j}(q_j)\ge v(q).
\end{equation}
This follows by applying lower semicontinuity to minimizing sources.
The function $v^\kappa(q)-|q|^2/(2\kappa)$ is concave, being an
infimum of affine functions.  Thus $v^\kappa$ is semiconcave and has
a full real second-order expansion almost everywhere.

Let $Q=(x,t,\xi_1,\ldots,\xi_m)$ denote the product coordinates,
understood as real coordinates when differentiating in time, and put
\[
 f_\kappa(Q)=\sum_i\alpha_i v^\kappa(x,\xi_i,t)-g_\theta(Q).
\]
This function is semiconcave.  The shared variables $x,t$ have not
been duplicated.  By \eqref{ell:epigraph-limit}, the positive boundary
gap on $O$ persists for all sufficiently small $\kappa$, whereas
$f_\kappa\le0$ at the original equality tuple.  For any sequence
$\delta_\kappa\downarrow0$, all minima of
$f_\kappa(Q)-p\cdot Q$ on $\overline O$, $|p|\le\delta_\kappa$,
are interior for small $\kappa$.  Any limit of such minimizing points
lies in the original equality set: the minimum values have upper
limit at most zero, and \eqref{ell:epigraph-limit} gives the reverse
inequality at each limit.

For completeness, these contacts form a set of positive Lebesgue
measure for each fixed small $\kappa$ and $\delta_\kappa>0$.
Let $E_\kappa$ be the set of minimizers for all slopes in that closed
ball.  Continuity of $f_\kappa$ and the boundary gap make $E_\kappa$
a compact subset of $O$.  At a contact, its supporting lower affine function and
semiconcavity imply differentiability, with $Df_\kappa=p$.
If $x,y\in E_\kappa$ are close and $C_\kappa$ is a local
semiconcavity constant, then
\[
 0\le f_\kappa(y)-f_\kappa(x)-Df_\kappa(x)\cdot(y-x)
 \le\tfrac12C_\kappa|x-y|^2.
\]
Combining the lower support at $x$ with the upper quadratic bound
based at $y$, evaluated at $y+h$, gives
\[
 (Df_\kappa(x)-Df_\kappa(y))\cdot h
 \le\tfrac12C_\kappa\bigl(|x-y|^2+|h|^2\bigr).
\]
Choose $h$ of length $|x-y|$ in the direction of the gradient
difference.  This proves that the gradient restricted to the contact
set is locally Lipschitz.  Its image contains the entire slope ball,
so that contact set cannot have measure zero.  Compact interior
localization makes the preceding argument valid in finitely many
coordinate balls.  Notice that no strictifying quadratic has been
added to $g_\theta$.

For each $i$, the exceptional set where $v^\kappa$ lacks an Alexandrov
expansion is null.  The projection
\[
 Q\longmapsto(x,t,\xi_i)
\]
is a surjective real linear map; its inverse image of that exceptional
set is null by Fubini.  A finite union remains null.  Hence we may
choose a contact $Q_\kappa\in E_\kappa$ at which every factor has a
full real second-order expansion simultaneously.  Write
\[
 a_i^\kappa=\partial_t v^\kappa(x_\kappa,\xi_i^\kappa,t_\kappa),
 \qquad W_i^\kappa=\dd_{x,\xi}v^\kappa
             (x_\kappa,\xi_i^\kappa,t_\kappa).
\]
Differentiating the common-base minimum gives
\begin{equation}\label{ell:time-sum}
 \sum_i\alpha_i a_i^\kappa
 =\partial_t\lambda(x_\kappa,t_\kappa)+(p_\kappa)_t,
 \qquad |p_\kappa|\le\delta_\kappa,
\end{equation}
and, for every complex $X,V_1,\ldots,V_m$,
\begin{equation}\label{ell:aggregate-levi}
 \sum_i\alpha_i W_i^\kappa[(X,V_i),\overline{(X,V_i)}]
 \ge\dd_x\lambda[X,\bar X]
      +(1-\theta)\dd\log\det A[V,\bar V].
\end{equation}
Here the forms on the right are evaluated at $Q_\kappa$.
To justify the passage from real to complex Hessians, if $H$ is a
real Hessian and $Z$ is a complex spatial direction, its Levi
quadratic form is
\[
 \tfrac14\bigl(H[Z_{\R},Z_{\R}]
                    +H[JZ_{\R},JZ_{\R}]\bigr).
\]
Apply the real Hessian inequality in both directions, with time
component zero.  The linear perturbation has zero Hessian, and
\eqref{ell:grassmann-penalty} gives
\eqref{ell:aggregate-levi}.  Thus this is an inequality of complex
Levi forms, rather than an identification of real and complex
Schur complements.

Let $C_i^\kappa$ denote the vertical block of $W_i^\kappa$.
Take $X=0$ and vary only the $i$th column in
\eqref{ell:aggregate-levi}.  Equation~\eqref{ell:single-column}
and convergence to the compact equality set give
\begin{equation}\label{ell:vertical-positive}
 C_i^\kappa\ge
 \bigl((1-\theta)(1-\max_j\alpha_j)-o(1)\bigr)I.
\end{equation}
The $o(1)$ tends to zero as $\kappa\downarrow0$, with the list and
$\theta$ fixed.  Replicate the list beforehand so that
$\max_i\alpha_i<1$; these vertical blocks are then uniformly positive.

We finally transfer these jets to the original envelope.  Fix
$\kappa$ and the chosen contact, and let $z_i$ be a minimizing source
in \eqref{ell:inf-convolution} for
$q_i=(x_\kappa,t_\kappa,\xi_i^\kappa)$.
Subtract $\varepsilon|h|^2$ from the real second-order expansion of
$v^\kappa(q_i+h)$.  For every $\varepsilon>0$ this gives a smooth
quadratic lower test $\phi_{i,\varepsilon}$ in a sufficiently small
neighborhood of $q_i$.  Since
\[
 v^\kappa(q_i+h)
 \le v(z_i+h)+\frac{|q_i-z_i|^2}{2\kappa},
\]
the translated function
\begin{equation}\label{ell:translated-test}
 \psi_{i,\varepsilon}(z_i+h)
 =\phi_{i,\varepsilon}(q_i+h)
       -\frac{|q_i-z_i|^2}{2\kappa}
\end{equation}
is a lower test of $v$ at $z_i$.  Translation preserves its time
derivative $a_i^\kappa$ and its Levi form
$W_i^\kappa-\varepsilon I$.

Apply Proposition~\ref{prop:disc-inequality}.  Its canonical
cotangent contraction has constant coefficients in these coordinates,
so the possibly different source base points and times cause no
coefficient error.  If $C_i^\kappa\ge cI$ and $0<\varepsilon<c/2$,
it supplies $K_{i,\varepsilon}\ge0$ with
\[
 K_{i,\varepsilon}\ge\Schur(W_i^\kappa-\varepsilon I),
 \qquad
 a_i^\kappa\ge
 \tr\bigl((C_i^\kappa-\varepsilon I)^{\mathsf T}
                  K_{i,\varepsilon}\bigr)
 \ge(c-\varepsilon)\tr K_{i,\varepsilon}.
\]
In particular $a_i^\kappa\ge0$.  At this fixed contact, let
$\varepsilon\downarrow0$ first.  The displayed trace bound gives a
convergent subsequence, and continuity of the Schur complement on
positive vertical blocks yields matrices $K_i^\kappa$ satisfying
\begin{equation}\label{ell:transferred-inequality}
 K_i^\kappa\ge0,\qquad
 K_i^\kappa\ge\Schur(W_i^\kappa),\qquad
 a_i^\kappa\ge\tr\bigl((C_i^\kappa)^{\mathsf T}K_i^\kappa\bigr).
\end{equation}
The temporary Hessian decrease has thus disappeared before taking
any limit in $\kappa$.  In particular it cannot be magnified by a
later vertical minimization.

\subsection{The scalar volume inequality}

\begin{proof}[Completion of the proof of Proposition~\ref{prop:ellipsoid-comparison}]
Use the preceding construction at the lower contact, with $B=I$.
Set $\mathcal K_\kappa=\sum_i\alpha_iK_i^\kappa$.
We first show
\begin{equation}\label{ell:schur-aggregate}
 \mathcal K_\kappa\ge0,
 \qquad \mathcal K_\kappa\ge\dd_x\lambda(x_\kappa,t_\kappa).
\end{equation}
For a Hermitian form with blocks
$W=\left(\begin{smallmatrix}H&D\\D^*&C\end{smallmatrix}\right)$,
$C>0$, completing the square gives
\[
 W[(X,V),\overline{(X,V)}]
 =X^*(H-DC^{-1}D^*)X
   +(V+C^{-1}D^*X)^*C(V+C^{-1}D^*X).
\]
Its minimum over the complex vertical vector $V$ is
$\Schur(W)[X,\bar X]$.  The extra term on the right of
\eqref{ell:aggregate-levi} is nonnegative by
\eqref{ell:rectangular-levi}.  Minimizing the left side independently
over every $V_i$ therefore proves
\[
 \sum_i\alpha_i\Schur(W_i^\kappa)
 \ge\dd_x\lambda(x_\kappa,t_\kappa).
\]
Equation~\eqref{ell:transferred-inequality} gives
\eqref{ell:schur-aggregate}.

Write $c_{\theta,\alpha}=(1-\theta)(1-\max_i\alpha_i)>0$.
Equations~\eqref{ell:time-sum}, \eqref{ell:vertical-positive}, and
\eqref{ell:transferred-inequality} imply
\begin{equation}\label{ell:aggregate-trace}
 \partial_t\lambda(x_\kappa,t_\kappa)+(p_\kappa)_t
 \ge\bigl(c_{\theta,\alpha}-o(1)\bigr)\tr\mathcal K_\kappa.
\end{equation}
The positive aggregate matrices are bounded.  Let $\kappa\downarrow0$
and pass to a subsequence to obtain a matrix $\mathcal K_{\theta,\alpha}$
such that
\[
 \mathcal K_{\theta,\alpha}\ge0,\qquad
 \mathcal K_{\theta,\alpha}\ge\dd_x\lambda(x_0,t_0),\qquad
 \partial_t\lambda(x_0,t_0)
       \ge c_{\theta,\alpha}\tr\mathcal K_{\theta,\alpha}.
\]
Now let $\theta\downarrow0$ and repeat the construction with increasingly
replicated lists, so that $\max_i\alpha_i\downarrow0$.
The aggregate matrices still act on the fixed $n$-dimensional
horizontal space and have uniformly bounded trace.  Another
subsequence gives
\begin{equation}\label{ell:limiting-majorant}
 \mathcal K\ge0,\qquad \mathcal K\ge\dd_x\lambda(x_0,t_0),
 \qquad \partial_t\lambda(x_0,t_0)\ge\tr\mathcal K.
\end{equation}
This also specifies the parameter order: fix the list and $\theta$;
regularize and select simultaneous contacts; remove the lower-test
Hessian decrease at each fixed contact; let $\kappa$ and the slopes
tend to zero; finally let $\theta$ and the maximum weight tend to zero.
No bound uniform in the number of individual contact jets is needed.

At the normalized point $B=I\le b_{t_0}$, and
$F_0=\log\det b_{t_0}$.  All eigenvalues of $b_{t_0}$ are at least
one, so each is at most their product.  Consequently
\[
 0<e^{-F_0}b_{t_0}\le I.
\]
Using the two matrix inequalities in
\eqref{ell:limiting-majorant}, in the indicated order, gives
\[
 \partial_t\lambda
 \ge\tr\mathcal K
 \ge e^{-F_0}\tr(b_{t_0}^{\mathsf T}\mathcal K)
 \ge e^{-F_0}\tr(b_{t_0}^{\mathsf T}\dd_x\lambda)
 =e^{-F_0}\Delta_{h_{t_0}}\lambda.
\]
The positivity of $\mathcal K$ justifies the middle comparison even
when $\dd_x\lambda$ is indefinite.  This proves
\eqref{ell:volume-viscosity}.
\end{proof}

\subsection{Equality of the envelope and the norm}

\begin{theorem}\label{thm:joint-positivity}
For the flow in Theorem~\ref{thm:localized-flow}, the function
\[
 (x,\xi,w)\longmapsto N(x,\xi,e^{w+\bar w})
\]
is plurisubharmonic on $T^{*(1,0)}M\times\C$.
For every finite $T$, its disc envelope satisfies $v=N$ on
$T^{*(1,0)}M\times(0,T)$.
\end{theorem}

\begin{proof}
In any holomorphic base coordinates let $\ell=\log\det b_t$.
The flow equation and the Ricci-form identity give
\[
 \partial_t\ell=R_{h_t},\qquad
 \dd_x\ell=\Ric(h_t),\qquad
 \Delta_{h_t}\ell=R_{h_t}.
\]
If $\varphi$ is an upper test of $F_0=\ell-\mathcal L$, then
$\lambda=\ell-\varphi$ is a lower test of $\mathcal L$.
Proposition~\ref{prop:ellipsoid-comparison} therefore says
\begin{equation}\label{ell:scalar-gap-equation}
 \partial_t\varphi
 \le e^{-F_0}\Delta_{h_t}\varphi
          +(1-e^{-F_0})R_{h_t}
\end{equation}
at contact.  The coefficient $F_0$ there equals the test value
$\varphi$.  By \eqref{ell:volume-gap}, $F_0$ is upper semicontinuous,
locally bounded, and between zero and $\rho$.  Since $\rho$ is smooth
and vanishes initially, extending $F_0$ by zero at $t=0$ preserves
upper semicontinuity.  Apply Proposition~\ref{prop:scalar-comparison}
to \eqref{ell:scalar-gap-equation} on each interval $[0,T']$ with
$T'<T$.  Since $T'$ is arbitrary, $F_0=0$ for every $t<T$.
Its proof does not require completeness of $h_t$ at positive time.

We have $B\le b_t$ and $\det B=\det b_t$.  The positive matrix
$b_t^{-1/2}Bb_t^{-1/2}\le I$ has determinant one; every one of its
eigenvalues is therefore one.  Hence $B=b_t$, and feasibility and
\eqref{ell:envelope-bounds} imply
\[
 N(x,\xi,t)=\xi^*B\xi\le v(x,\xi,t)\le N(x,\xi,t).
\]
This proves $v=N$.

For the raw disc infimum $v^{\rm raw}$ of \eqref{var:envelope},
we have $v\le v^{\rm raw}$, while constant discs give $v^{\rm raw}\le N$.
Thus $v^{\rm raw}=N$ too, so $N$ obeys the submean inequality for
every admissible disc, not merely for regularized centers.
Given a holomorphic disc $(f,W)$ in $Y\times\C$, put
$t_0=e^{W(0)+\overline{W(0)}}$ and $w=W-W(0)$.
Choose $T$ larger than the maximum of $e^{W+\bar W}$ on the closed
disc.  The defining disc inequality then gives
\[
 N(f(0),t_0)
 \le\frac1{2\pi}\int_0^{2\pi}
 N\bigl(f(e^{i\vartheta}),e^{W(e^{i\vartheta})+
                  \overline{W(e^{i\vartheta})}}\bigr)\,d\vartheta.
\]
It suffices to use discs smooth past their boundary, including all
sufficiently small coordinate discs.  Since the joint norm is smooth,
these submean inequalities are precisely its plurisubharmonicity.
The horizon $T$ was arbitrary.
\end{proof}
